\documentclass[10pt,twocolumn]{article}
\usepackage[T1]{fontenc}
\usepackage[utf8]{inputenc}
\usepackage[brazil]{babel}
\usepackage{lmodern,booktabs,graphicx,microtype,geometry}
\usepackage[hidelinks]{hyperref}
\usepackage{tikz}
\usetikzlibrary{arrows.meta,positioning,fit,backgrounds,shapes.geometric}
\geometry{a4paper,margin=1.9cm}

\usepackage{titlesec}
\titlespacing*{\section}{0pt}{1.1em}{0.4em}
\titlespacing*{\subsection}{0pt}{0.8em}{0.3em}
\setlength{\parskip}{0pt}

\hyphenation{har-ness su-per-cé-re-bro au-di-ta-bi-li-da-de ras-tre-á-vel
  me-to-do-lo-gia con-fir-ma-ção per-mis-sões}

\newcommand{\code}[1]{\texttt{\footnotesize #1}}

\title{\textbf{Lastro}: um harness de grafo com ReAct nos nós e permissão
\emph{deny-first} para gestão de marketing}
\author{Davi Eduardo}
\date{Agosto de 2026}

\begin{document}
\maketitle

\begin{abstract}
\noindent
O gestor de conta de uma agência não faz pedidos de um passo. ``Por que caíram as
vendas da Ômega 3 essa semana?'' exige resolver uma entidade na memória da conta,
consultar duas ou três APIs, cruzar identificadores, formular e descartar hipóteses ---
e, no fim, defender o número numa reunião com o cliente. Este paper propõe
\textbf{Lastro}, um harness híbrido para o chat agêntico da AdzHub:
\textbf{grafo de estados como espinha dorsal}, \textbf{loops ReAct dentro dos nós de
investigação} e uma \textbf{camada de permissões deny-first}, com o supercérebro da conta
--- grafo de entidades mais linha do tempo --- tratado como fonte de contexto de primeira
classe, acessada por tools dedicadas em vez de despejada no prompt. O argumento central do
estudo é de enquadramento: os cinco tipos de harness de referência do desafio não estão no
mesmo nível de abstração. ReAct é um loop; grafo é uma topologia; sandbox é um ambiente de
execução; permissões e skills são uma camada transversal; RLM é uma estratégia de acesso a
contexto. Só os dois primeiros disputam o mesmo lugar, e a pergunta certa não é ``qual
escolher'', é ``o que fica por fora, o que fica por dentro e o que atravessa tudo''. A
aposta de produto que sustenta o desenho: em marketing, a \emph{defensibilidade} da
resposta vale mais que a autonomia do agente. Dois recortes são declarados com o custo de
cada um: sem sandbox/CodeAct e sem RLM completo --- apenas compactação por nó. O paper
nomeia onde a arquitetura quebra nas quatro tarefas reais do gestor, em vez de alegar
cobertura. O desenho não ficou no papel: o harness roda numa demo pública, onde os quatro
percursos do gestor executam ponta a ponta e o gate se comporta como a arquitetura exige ---
aprovar executa a ação uma vez, reenviar a mesma aprovação executa zero.
\end{abstract}

\smallskip
\noindent\textbf{Palavras-chave:} harness agêntico; grafo de estados; ReAct;
permissões deny-first; memória temporal; gestão de marketing.

% ---------------------------------------------------------------------------
\section{Introdução}

O produto-alvo da AdzHub é um chat agêntico para o profissional de marketing: o gestor
descreve uma tarefa em linguagem natural e o sistema orquestra contexto e ferramentas até
produzir algo útil na operação. Abaixo do chat existem três camadas: o \emph{supercérebro}
(memória da conta, um grafo de pessoas, campanhas, canais e decisões somado a uma linha do
tempo), os \emph{Apps} (metodologias da agência empacotadas: insight da semana, brief de
criativo, diagnóstico de conta, análise de criativos) e as \emph{APIs} da operação (Meta
Ads, Google Ads, Analytics, CRM, WhatsApp). A conta de referência é a Housewhey,
e-commerce de suplementos atendido pela SPOT, com o time de Aline, Carolina e Luiza.

Convém separar duas coisas que a conversa pública mistura. O \textbf{modelo} é a função
que, dado um contexto, escolhe a próxima ação. O \textbf{harness} é o runtime em volta: ele
decide quais ações existem, o que o modelo enxerga antes de escolher, quantas vezes pode
tentar, o que é feito com cada resultado, o que fica registrado e o que exige um humano
antes de acontecer. A documentação do AI SDK decompõe um agente exatamente nesses três
elementos --- LLM, tools e loop --- e é o terceiro que este paper trata~\cite{vercel2026aisdk}.
Trocar o modelo melhora a média das respostas. Trocar o harness muda o que o sistema
consegue prometer.

\paragraph{De chatbot a agente.} Três propriedades fazem a diferença no dia a dia do gestor,
e nenhuma delas é ``responder melhor''. Primeira: o sistema \textbf{resolve entidade contra
a operação real}. ``A Ômega 3'' vira um \code{id} de campanha da Housewhey e uma janela de
datas concreta, antes de qualquer chamada de API. Segunda: ele \textbf{busca o dado que ele
mesmo decidiu que faltava}. Num diagnóstico de queda de vendas, ninguém programou ``olhar o
link de destino dos criativos''; esse passo nasce da observação de que o Meta reporta mais
conversão do que o CRM atribui. Terceira: ele \textbf{age sob permissão} --- pausa anúncio
de verdade, depois de mostrar o que vai pausar e quanto aquilo gastou. Chatbot devolve
texto; agente devolve efeito rastreável.

\paragraph{Tese.} Proponho um harness híbrido: grafo de estados como espinha dorsal, loops
ReAct dentro dos nós de investigação e permissões deny-first como camada transversal, com o
supercérebro acessado por tool. O grafo dá fronteiras nomeadas --- \code{interpret},
\code{plan}, \code{fetch}, \code{reason}, \code{gate}, \code{act}, \code{respond} --- que
tornam um diagnóstico de cinco etapas legível para quem vai repeti-lo numa call com o
cliente. O ReAct dentro do nó preserva a liberdade de descobrir o passo que ninguém
planejou, sob teto de passos e allowlist própria. O gate garante que nenhuma ação com
efeito real acontece sem um humano ver o preview do efeito. O nome do harness resume a
aposta: nenhuma afirmação circula sem o lastro que a segura.

\paragraph{Mapa.} A Seção~2 fixa o vocabulário e lê o domínio. A Seção~3 é o coração
argumentativo: por que grafo mais ReAct mais permissão, por que não cada uma das outras
referências isoladamente, o que é tool, o que é memória, o que é App, o que mudou no meu
entendimento durante o estudo e quais trade-offs foram aceitos de propósito. A Seção~4
descreve o harness que foi construído e está no ar, com um turno percorrido nó por nó, e a
Seção~5 diz como executá-lo --- chave, modelo, dados e o que os testes cobrem. A Seção~6 nomeia onde a
arquitetura quebra nas quatro tarefas reais do gestor, e a Seção~7 diz o que eu construiria
em seguida --- e o que deliberadamente não construiria.

% ---------------------------------------------------------------------------
\begin{figure*}[t]
\centering
\resizebox{\textwidth}{!}{%
\begin{tikzpicture}[
  font=\small,
  b/.style={draw, rounded corners=2pt, align=center, inner sep=4pt,
            text width=30mm, minimum height=11mm},
  hb/.style={b, very thick, text width=34mm, minimum height=15mm},
  gt/.style={draw, trapezium, trapezium left angle=72, trapezium right angle=108,
             align=center, inner sep=3pt, text width=26mm},
  fl/.style={-{Latex[length=2mm]}, thick},
  rt/.style={-{Latex[length=2mm]}, dashed},
  lb/.style={font=\scriptsize, inner sep=1.5pt}
]
\node[b]  (g)   at (0,0)        {Gestor de marketing\\[1pt]{\scriptsize Aline $\cdot$ Carolina $\cdot$ Luiza}};
\node[b]  (c)   at (4.4,0)      {Chat agêntico};
\node[hb] (h)   at (9.0,0)      {\textbf{HARNESS}\\[1pt]{\scriptsize grafo de estados\\ ReAct nos nós\\ permissões \emph{deny-first}}};
\node[b]  (sc)  at (14.6,2.6)   {Supercérebro\\[1pt]{\scriptsize grafo + linha do tempo}\\[1pt]\textsc{\scriptsize memória}};
\node[b]  (ap)  at (14.6,0)     {Apps de metodologia\\[1pt]{\scriptsize diagnóstico, brief, criativos}\\[1pt]\textsc{\scriptsize app}};
\node[b]  (api) at (14.6,-2.6)  {APIs da operação\\[1pt]{\scriptsize Meta, Google, GA, CRM, WhatsApp}\\[1pt]\textsc{\scriptsize api}};
\node[gt] (gate)at (9.0,-4.2)   {\textbf{gate}\\[1pt]{\scriptsize confirmação humana}};
\node[b]  (act) at (14.6,-4.6)  {\textbf{act}\\[1pt]{\scriptsize ação com efeito real}};

\draw[fl] (g)  -- node[lb,above]{pedido} (c);
\draw[fl] (c)  -- node[lb,above]{turno} (h);
\draw[rt] (h.north) to[out=120,in=60] node[lb,above]{resposta, artefatos e trace} (c.north);
\draw[rt] (c) to[bend left=22] node[lb,below]{} (g);

\draw[fl] (h.east) to[bend left=11] (sc.west);
\draw[rt] (sc.west) to[bend left=11] (h.east);
\draw[fl] (h.east) to[bend left=4] node[lb,above,pos=0.55]{tool de leitura} (ap.west);
\draw[rt] (ap.west) to[bend left=4] (h.east);
\draw[fl] (h.east) to[bend right=11] (api.west);
\draw[rt] (api.west) to[bend right=11] (h.east);

\draw[fl] (h.south) -- node[lb,right]{tool de escrita (proposta)} (gate.north);
\draw[fl] (gate.east) -- node[lb,above]{confirmado} (act.west);
\draw[fl] (act.north) -- (api.south);
\draw[rt] (gate.north west) to[out=150,in=210] node[lb,left]{negado} (h.south west);
\end{tikzpicture}}
\caption{O harness entre o gestor e as três camadas. O \textbf{supercérebro} é
\textbf{memória} (grafo de entidades mais linha do tempo); os \textbf{Apps} são
\textbf{metodologia} empacotada, que devolve estrutura e não parágrafo; as \textbf{APIs} são
o dado bruto da operação. As três são acessadas por um mecanismo único --- \textbf{tool} ---
e nada entra no contexto do modelo por outro caminho. Setas cheias são chamadas do harness
para fora; tracejadas são retornos que viram observação no estado. O \textbf{gate não é
tool}: é um nó do grafo, que interrompe o turno e devolve um preview em português antes de
qualquer escrita.}
\label{fig:sistema}
\end{figure*}

\section{Preliminares}

\subsection{Harness, no vocabulário deste paper}

\textbf{Loop.} A regra de repetição: o modelo propõe, o runtime executa, o resultado
volta e o ciclo recomeça até uma condição de parada. \textbf{Tool} é uma capacidade
declarada com assinatura e argumentos validados; adoto o contrato
\emph{Action--Execution--Observation} do OpenHands SDK, em que o LLM propõe um JSON, o
runtime valida contra um esquema, um executor roda e o resultado retorna como
\emph{observação}~\cite{wang2026openhands}. \textbf{Observação} é justamente esse retorno já
convertido em algo que entra no contexto --- e, neste desenho, é a \emph{única} coisa que
entra. \textbf{Allowlist} é o conjunto de tools visível num dado momento; o AI SDK entrega
isso por \code{prepareStep} com \code{activeTools}, sem grafo
nenhum~\cite{vercel2026loopcontrol}. \textbf{Teto de passos} é o freio de custo: na versão
7.x do AI SDK, \code{stopWhen} com \code{isStepCount(20)} é o
padrão\footnote{\code{maxSteps} não aparece mais na documentação 7.x; o controle atual é
\code{stopWhen}. Registro a correção porque material de 2025 ainda cita o nome
antigo.}~\cite{vercel2026loopcontrol}. \textbf{Estado} é a estrutura tipada que sobrevive
entre etapas --- no LangGraph, um \code{State} em que cada chave pode ter um \emph{reducer}
dizendo se o update substitui ou acumula~\cite{langchain2026langgraph}. \textbf{Checkpoint}
é o ponto onde esse estado é salvo; no LangGraph, nas fronteiras de \emph{superstep} e
nunca no meio de um nó. \textbf{Gate} é o ponto em que o turno para e devolve a decisão a um
humano.

Vale dizer o que \emph{não} é harness: prompt de sistema não é harness, e RAG no prompt
tampouco. Ambos influenciam o modelo sem criar fronteira verificável nenhuma.

\subsection{O domínio, como eu o leio}

As três camadas do enunciado não são três fontes de dado equivalentes. O supercérebro é
\textbf{memória}: responde ``quem é quem'' e ``o que aconteceu quando''. Os Apps são
\textbf{metodologia}: são o ativo da agência, o jeito da SPOT de diagnosticar uma conta ou
escrever um brief. As APIs são \textbf{dado bruto} da operação. A Figura~\ref{fig:sistema}
mostra como o harness se põe entre o gestor e as três.

Mais importante: as quatro tarefas típicas do gestor \textbf{não são a mesma classe de
problema}. T1, o relatório de criativos cruzado com resultado real por \code{utm\_content},
é um \emph{join determinístico} que precisa dar o mesmo número toda semana. T2, o
diagnóstico de anomalia, é uma \emph{investigação aberta} cujo caminho ninguém enumera
antes. T3, a pauta de reunião, é \emph{navegação de um histórico} que só cresce. T4, a
análise de criativos com proposta de pausa, \emph{termina numa ação irreversível} sobre a
verba de um cliente. Um harness sintonizado para uma delas está desafinado para as outras
três --- e é daí que o híbrido nasce, não de indecisão.

% ---------------------------------------------------------------------------
\begin{table}[t]
\centering
\footnotesize
\caption{Catálogo de tools desenhado para o domínio: 10 de leitura, 2 de escrita, toda
escrita atravessando o \code{gate}. A tabela é \textbf{ilustrativa} do contrato de
fronteira --- não é uma allowlist de produção, que seria por conta, por perfil de usuário e
versionada.}
\label{tab:tools}
\begin{tabular}{@{}llll@{}}
\toprule
Tool & Camada & Efeito & Nó \\
\midrule
\code{graph\_query}        & supercérebro & read  & int., fet., rea. \\
\code{timeline\_query}     & supercérebro & read  & int., fet., rea. \\
\code{meta\_ads\_insights} & API          & read  & fetch \\
\code{google\_ads\_insights}& API         & read  & fetch \\
\code{ga\_report}          & API          & read  & fetch \\
\code{crm\_leads}          & API          & read  & fetch \\
\code{list\_criativos}     & API          & read  & fetch \\
\code{get\_metrics}        & API          & read  & fetch \\
\code{app\_diagnostico}    & App          & read  & reason \\
\code{propose\_ctas}       & App          & read  & reason \\
\code{pause\_ads}          & API          & \textbf{write} & act (pós-gate) \\
\code{send\_whatsapp}      & API          & \textbf{write} & act (pós-gate) \\
\bottomrule
\end{tabular}
\end{table}

\section{Projeto}

\subsection{Cinco referências, três níveis de abstração}

O enunciado lista cinco tipos de harness e pede que se escolha um, combine ou invente. A
primeira coisa que o estudo produziu foi a suspeita de que a lista é uma armadilha de
enquadramento --- não por erro do enunciado, mas porque os cinco itens não descrevem o mesmo
tipo de objeto:

\begin{itemize}\itemsep2pt
\item \textbf{ReAct} é um \emph{loop}: uma regra sobre como intercalar raciocínio e
ação~\cite{yao2023react}.
\item \textbf{Grafo de estados} é uma \emph{topologia}: quais etapas existem e como se
ligam~\cite{langchain2026langgraph}.
\item \textbf{Sandbox/CodeAct} é um \emph{ambiente de execução}: onde a ação roda e com que
isolamento~\cite{wang2024codeact}.
\item \textbf{Permissões e skills} é uma \emph{camada transversal}: o que é permitido e qual
instrução está carregada~\cite{anthropic2026permissions}.
\item \textbf{RLM} é uma \emph{estratégia de acesso a contexto}: o contexto é injetado ou
navegado~\cite{zhang2025rlm}.
\end{itemize}

Só os dois primeiros disputam de fato o mesmo lugar. Os outros três compõem com qualquer
coisa. A pergunta de arquitetura, portanto, não é ``qual dos cinco''; são três perguntas
separadas: \textbf{o que fica por fora} (a topologia), \textbf{o que fica por dentro} (o
loop) e \textbf{o que atravessa tudo} (permissão, contexto, ambiente). Este paper responde:
grafo por fora, ReAct por dentro, permissão deny-first atravessando; ambiente de execução
recusado; contexto acessado por consulta.

A Tabela~\ref{tab:tipos} resume o julgamento a que cheguei sobre os cinco tipos
\emph{puros}, aplicados a este domínio. Ela torna difícil negar três coisas.
\textbf{(i)} Auditabilidade e flexibilidade são anticorrelacionadas nos tipos puros --- os
dois que pontuam alto em flexibilidade pontuam baixo em auditabilidade, e vice-versa; não é
coincidência, já que auditabilidade vem de estrutura declarada \emph{antes} e flexibilidade
vem de estrutura decidida \emph{durante}. \textbf{(ii)} A coluna de segurança operacional é
quase ortogonal às outras, o que reforça que permissões não são um tipo concorrente e sim
uma camada. \textbf{(iii)} Ninguém combina custo baixo com flexibilidade alta exceto o
ReAct, que paga a conta na auditabilidade.

\begin{table}[t]
\centering
\footnotesize
\scriptsize
\caption{Os cinco tipos \emph{puros} avaliados para o domínio da AdzHub. Notas de 1 (ruim)
a 5 (ótimo). Colunas: auditabilidade, flexibilidade, segurança operacional, custo/latência
e esforço de modelagem. \textbf{É julgamento do autor} a partir das mecânicas descritas nas
fontes, não medição de nenhuma delas.}
\label{tab:tipos}
\setlength{\tabcolsep}{4pt}
\begin{tabular}{@{}lccccc@{}}
\toprule
Tipo & Aud. & Flex. & Seg. & Custo & Model. \\
\midrule
ReAct                  & 2 & 5 & 2 & 4 & 5 \\
CodeAct / sandbox      & 2 & 5 & 1 & 3 & 2 \\
Permissões \& skills   & 3 & 4 & 5 & 4 & 3 \\
Grafo de estados       & 5 & 2 & 4 & 4 & 1 \\
RLM / ambiente         & 1 & 5 & 2 & 1 & 3 \\
\bottomrule
\end{tabular}
\end{table}

\subsection{Por que nenhuma delas sozinha}

Cada tipo puro foi executado mentalmente contra uma tarefa real da Housewhey até quebrar, e
cada falha recebeu nome. O resumo:

\textbf{ReAct puro --- a conclusão órfã.} No diagnóstico T2 o loop roda doze ou treze
chamadas e escreve um parágrafo correto. Na terça-feira, na call, a Aline pergunta ``de onde
saiu que a taxa de comparecimento caiu?''. O número existe na resposta e não existe em lugar
nenhum do traço como resultado nomeado: foi calculado de cabeça a partir de dois JSONs que
passaram pelo contexto seis passos antes. Para auditar, o gestor teria que ler as treze
observações cruas --- exatamente o trabalho que ele delegou. Pior: numa conversa longa entra
a compactação. O OpenHands documenta que seu \code{Condenser} descarta eventos e os
substitui por sumários, e que é isso que reduz custo de API em até
2$\times$~\cite{wang2026openhands}; num ReAct puro, sem log fora da janela, a observação que
sustentava a conclusão pode ter sido resumida para fora antes da resposta final. A resposta
sobrevive; a evidência, não. A interpretabilidade que o ReAct reivindica é a \emph{do
traço}, legível por quem lê o traço inteiro --- não é a auditabilidade estrutural de que o
gestor precisa.

\textbf{CodeAct puro --- a metodologia volátil.} T1 é onde o CodeAct deveria ganhar de
lavada, e quase ganha: cruzar gasto por anúncio com leads do CRM é literalmente a
``composição de múltiplas tools'' que o paper original diz que o formato JSON não
faz~\cite{wang2024codeact}. Duas falhas aparecem depois. A primeira: a definição da SPOT do
que conta como lead qualificado passa a viver no código que o modelo escreveu \emph{naquele}
turno; na semana seguinte ele escreve outro, trata os leads órfãos de outro jeito, e o
cliente recebe dois relatórios diferentes sem que nada tenha mudado na conta. O App de
metodologia deixa de ser App e vira uma sugestão que o modelo segue quando lembra --- para
uma empresa cujo produto \emph{é} a metodologia empacotada, trocar App por script gerado é
vender o ativo. A segunda: para fazer o join, os leads entram no processo com nome, telefone
e e-mail. O OpenHands trata a versão \emph{credencial} desse problema com um
\code{SecretRegistry} que injeta segredos só na execução e mascara ocorrências na
saída~\cite{wang2026openhands} --- mas isso protege o token do Meta Ads, não o telefone do
lead. Um sandbox de propósito geral não sabe o que é dado pessoal.

\textbf{Permissões e skills sozinhas --- a ordem não governada.} O fluxo de avaliação do
Claude Agent SDK (hooks $\rightarrow$ deny $\rightarrow$ ask $\rightarrow$ modo
$\rightarrow$ allow $\rightarrow$ \emph{callback}) é a melhor especificação de segurança
disponível de graça~\cite{anthropic2026permissions}, e as \emph{skills} do mesmo SDK
resolvem elegantemente o problema dos Apps de metodologia~\cite{anthropic2026agentsdk}. Mas permissão governa \emph{o quê}, nunca \emph{quando}.
Em T4, nada no sistema de permissões tem opinião sobre sequência: o agente pode listar os
criativos, ler as copies e escrever três variações de CTA \emph{antes} de olhar o
desempenho, porque escrever copy é a parte que ele faz bem e a métrica é a parte chata. O
resultado é um briefing bonito derivado do texto do criativo em vez do dado. Não existe
regra de permissão que expresse ``não proponha CTA antes de ter lido a métrica'', porque
isso não é uma pergunta sobre autorização.

\textbf{Grafo puro --- a enumeração antecipada.} Um grafo fixo
(\code{coletar} $\rightarrow$ \code{pendências} $\rightarrow$ \code{montar\_pauta}) resolve
T3 muito bem, toda segunda-feira, de forma auditável. Aí o gestor lê e digita: ``e o que
ficou pendente com a Luiza no WhatsApp semana passada?''. Não há nó para isso. As opções são
todas ruins: alucinar com o que está no estado, responder que não sabe fazer, ou abrir o
código e acrescentar um nó --- que é a resposta certa em engenharia e a errada num produto
de chat, onde a pergunta seguinte é sempre imprevisível. O gestor de marketing não conversa
em fluxograma, e o próprio enunciado diz que não predefine o que o chat precisa fazer.

\textbf{RLM puro --- latência sem teto e auditoria descartável.} O princípio é o certo: em
vez de injetar o contexto, deixar o modelo navegá-lo programaticamente. Os ganhos são
grandes onde o corpus não cabe na janela --- em OOLONG a 132k tokens, o RLM sobre GPT-5-mini
atinge cerca de 64 pontos contra cerca de 30 do modelo puro, com custo por query
aproximadamente igual~\cite{zhang2025rlm}. Mas os próprios autores declaram que as chamadas
recursivas são bloqueantes, que a duração vai de segundos a ``vários minutos'', que não há
garantia forte sobre custo total, e que a degradação é maior justamente em problemas de
\emph{contagem}. Um chat interativo tem um orçamento de paciência de dezenas de segundos, e
o pior caso não é o caso médio: é o caso que acontece na frente do cliente. Some-se a isso
que a procedência de um ``3 decisões pendentes'' passa a ser um trecho de Python que o
modelo escreveu, executou e não guardou.

\subsection{A escolha: grafo por fora, ReAct por dentro, permissão atravessando}

\paragraph{Por fora, o grafo.} O harness é um grafo de estados com nós nomeados e edges
condicionais nomeadas (Figura~\ref{fig:grafo}). O estado é explícito e serializável; cada
transição é um checkpoint. Três coisas passam a existir por construção. Primeira,
\textbf{allowlist por fase}: \code{pause\_ads} só existe no nó \code{act} --- não é uma
instrução de prompt que o modelo pode ignorar, é topologia. Segunda, \textbf{trace legível
como narrativa}: eventos de entrada e saída de nó dão à interface a legenda
\emph{pedido $\rightarrow$ raciocínio $\rightarrow$ tool $\rightarrow$ observação
$\rightarrow$ ação $\rightarrow$ resposta} sem que ela precise inferir nada. Terceira,
\textbf{retomada}: o gate interrompe o turno de verdade, e voltar é reentrar num nó com o
estado restaurado, não recomeçar a conversa.

A objeção padrão ao grafo é que auditoria custa caro. Ela não se sustenta com números na
mesa. O OpenHands mede o overhead do seu event sourcing em 39.870 eventos de 433 conversas:
\textbf{0,20\,ms de persistência por evento} (mediana; 0,31\,ms no P95), 4,1\,ms para
reconstruir o estado inteiro por replay e 380\,KB de armazenamento por conversa
(mediana)~\cite{wang2026openhands}. No mesmo trabalho, o redesenho arquitetural centrado
nesse log acompanhou uma \textbf{queda de 61\% nos erros atribuíveis ao sistema} --- de 78,0
para 30,0 por mil conversas --- num rollout de produção de 15 dias com as duas versões
servindo usuários em paralelo. O custo de um agente é o LLM, não o log. Quem economiza
estrutura não está economizando dinheiro; está economizando a única coisa que torna a
resposta defensável.

\paragraph{Por dentro, ReAct.} Se cada nó fizesse uma chamada fixa de tool, o harness seria
um workflow determinístico --- e workflow determinístico não diagnostica. Dentro de
\code{fetch} e \code{reason} roda um loop ReAct: pensar, chamar tool, observar, repetir,
limitado por três coisas --- teto de passos por nó, allowlist do nó (checada duas vezes, na
configuração do nó \emph{e} na definição da tool) e um teto de tokens de observação que,
estourado, desvia para o nó \code{compact} em vez de truncar. O ciclo \code{fetch}
$\leftrightarrow$ \code{reason} tem teto próprio; ao estourar, o grafo força a saída para
\code{respond}, que redige \emph{declarando a lacuna} em vez de preencher com
plausibilidade. É essa liberdade interna que permite o passo que ninguém planejou: buscar o
link de destino dos criativos só faz sentido \emph{depois} de observar que o Meta reporta
mais conversão do que o CRM atribui.

\paragraph{Atravessando, a permissão.} Toda tool declara efeito \code{read} ou
\code{write}. Leitura roda livre dentro da allowlist. \emph{Toda} escrita para o turno no nó
\code{gate}, que monta um preview em português --- itens afetados pelo nome que o gestor
reconhece e com o número que justifica; impacto em uma frase; se é reversível e como
desfazer; e o que acontece se ele negar --- e devolve o controle. O desenho é deny-first
porque a fonte mais explícita sobre isso avisa que uma checagem colocada tarde demais é
silenciosamente contornada: no Claude Agent SDK, \emph{tools auto-aprovadas nunca chegam ao
callback} \code{canUseTool}, e por isso a verificação que precisa valer para toda chamada
tem que morar num hook que roda antes de todo o resto~\cite{anthropic2026permissions}. A
armadilha correspondente aqui é fácil de imaginar: alguém marca a leitura do Meta Ads como
permitida, alguém acrescenta a pausa de anúncios ao mesmo conector, e a confirmação some. A
separação do OpenHands entre \emph{quem avalia risco} e \emph{quem aplica a
política}~\cite{wang2026openhands} é o que permite política diferente por conta sem
reescrever tool.

Vale registrar que esse vocabulário é o da própria casa: a descrição do episódio do podcast
da AdzHub sobre harness anuncia o loop ReAct, a compressão de contexto e as políticas
deny-first entre seus temas~\cite{adzhub2026harness}.\footnote{Não obtive transcrição nem
áudio do episódio; a única coisa acessível foi a descrição publicada. Cito como descrição de
episódio, não como argumento técnico. O mesmo vale para o episódio do \emph{How I
AI}~\cite{vo2026harness}, do qual li apenas as notas.}

\begin{figure*}[t]
\centering
\resizebox{\textwidth}{!}{%
\begin{tikzpicture}[
  font=\normalsize,
  n/.style={draw, rounded corners=2pt, align=center, inner sep=3pt,
            text width=22mm, minimum height=10mm},
  t/.style={draw, rounded corners=7pt, align=center, inner sep=3pt,
            text width=20mm, minimum height=9mm, densely dotted},
  gt/.style={draw, diamond, aspect=2.6, align=center, inner sep=1pt},
  fl/.style={-{Latex[length=2.4mm]}, thick},
  er/.style={-{Latex[length=2.4mm]}, dashed},
  lb/.style={font=\small, inner sep=1.5pt}
]
\node[n] (int)   at (0,0)     {\code{interpret}\\[1pt]{\small resolve entidades}};
\node[n] (plan)  at (3.4,0)   {\code{plan}\\[1pt]{\small monta os passos}};
\node[n] (fet)   at (7.0,0)   {\code{fetch}\\[1pt]{\small tools de leitura}};
\node[n] (rea)   at (10.6,0)  {\code{reason}\\[1pt]{\small cruza e descarta}};
\node[gt](gate)  at (14.4,0)  {\code{gate}\\[1pt]{\small efeito real?}};
\node[n] (act)   at (17.6,0)  {\code{act}\\[1pt]{\small executa o confirmado}};
\node[n] (resp)  at (21.0,0)  {\code{respond}\\[1pt]{\small redige + artefatos}};

\node[n] (err)   at (4.4,-3.4)  {\code{errorHandler}\\[1pt]{\small retry ou degrada}};
\node[n] (comp)  at (10.6,-3.4) {\code{compact}\\[1pt]{\small resume observações}};
\node[t] (halt)  at (14.4,-3.4) {\textbf{turno pausa}\\[1pt]{\small aguardando humano}};

\draw[fl] (-2.0,0) -- node[lb,above=3pt]{turno inicia} (int);
\draw[fl] (int)  -- node[lb,above=20pt]{entidades resolvidas} (plan);
\draw[fl] (plan) -- node[lb,above=20pt]{precisa dados} (fet);
\draw[fl] (fet)  -- node[lb,below=20pt]{dados coletados} (rea);
\draw[fl] (rea)  to[bend right=22] node[lb,above=2pt]{lacuna de dado} (fet);
\draw[fl] (rea)  -- node[lb,below=20pt]{conclusão pede ação} (gate);
\draw[fl] (act)  -- (resp);
\draw[fl] (resp) -- node[lb,above=3pt]{turno encerra} (23.0,0);

\draw[fl] (rea.south west) -- node[lb,left,pos=0.55]{contexto estourou} (comp.north west);
\draw[fl] (comp.north east) -- node[lb,right,pos=0.45]{compactado} (rea.south east);

\draw[fl] (gate.south) -- node[lb,right]{aguardando confirmação} (halt.north);
\draw[fl] (halt.east) to[out=0,in=250] node[lb,right,pos=0.55]{confirmada} (act.south);
\draw[fl] (halt.south) to[out=-25,in=250] node[lb,below,pos=0.45]{negada} (resp.south);
\draw[fl] (gate.north) to[out=60,in=120] node[lb,above]{sem efeito real} (resp.north);
\draw[fl] (rea.north) to[out=62,in=118] node[lb,above]{conclusão sem ação} (resp.north);

\draw[er] (fet.south) to[out=-90,in=90] node[lb,right,pos=0.5]{falha de tool} (err.north);
\draw[er] (err.north west) to[out=90,in=200] node[lb,left,pos=0.5]{retry} (fet.west);
\draw[er] (err.south) to[out=-30,in=250] node[lb,below,pos=0.35]{degradar} (resp.south west);

\begin{scope}[on background layer]
\node[draw, dashed, rounded corners=4pt, inner sep=13pt,
      fit=(fet)(rea),
      label={[font=\normalsize]above:\textbf{loop ReAct} --- pensar, chamar tool, observar, repetir; teto de passos e allowlist por nó}] {};
\end{scope}
\end{tikzpicture}}
\caption{O grafo de estados interno. O \textbf{loop ReAct vive apenas dentro de
\code{fetch} e \code{reason}} (caixa tracejada): lá o modelo escolhe tool, observa e decide o
próximo passo, sob teto de passos e allowlist própria. Fora dela não há loop --- os demais
nós são passagem única. O \textbf{gate interrompe o turno de verdade}: a resposta HTTP fecha
em ``aguardando confirmação'' e a execução só retoma num novo pedido com a decisão do
gestor; não existe caminho de \code{reason} direto para \code{act}. A aresta de \code{act}
para \code{respond}, sem rótulo no desenho por falta de espaço, chama-se
\code{acao\_confirmada} no runtime. Omitidas para
legibilidade: as saídas diretas de \code{interpret} (ambiguidade de entidade: o agente
pergunta em vez de chutar) e de \code{plan} (pedido respondível sem dado novo) para
\code{respond}, e o caminho de erro que sai de \code{act}, idêntico ao de \code{fetch}.}
\label{fig:grafo}
\end{figure*}

\subsection{Gate e efeito precisam ser nós separados}

Esta é a consequência de projeto mais cara que o estudo produziu, e ela vem de uma linha da
documentação do LangGraph que a maioria dos tutoriais ignora: efeitos colaterais anteriores
a um \code{interrupt()} \emph{devem ser idempotentes}, porque, na retomada, \textbf{o nó
afetado roda de novo desde o início da função}~\cite{langchain2026langgraph}. O checkpoint
salva nas fronteiras de superstep, não no meio do nó.

O que isso significa em produção de marketing é direto. Se o mesmo nó pede a aprovação
\emph{e} dispara o WhatsApp para a Aline, a retomada manda a mensagem duas vezes. Não existe
``despausar'' uma mensagem enviada ao cliente. Daí a regra não-negociável deste desenho: um
nó \code{gate} que \emph{só} interrompe e monta o preview, e um nó \code{act} que \emph{só}
executa --- e que executa exatamente os argumentos que estavam no preview, sem re-inferir
nada, para que o gestor não aprove uma coisa e outra aconteça. Não existe caminho de
\code{reason} direto para \code{act}.

\subsection{O que é tool, o que é memória, o que é App}

A regra que organiza a fronteira inteira cabe em uma frase: \textbf{tudo que entra no
contexto do modelo entra como retorno de tool}. Não existe contexto ambiente, não existe
bloco de conta colado no prompt de sistema. Memória, App e API são \emph{camadas}; tool é o
\emph{mecanismo único de acesso} a todas elas. A consequência prática é que cada afirmação
da resposta tem um evento de trace com uma fonte citável atrás dela.

\textbf{Memória} é o supercérebro, acessado por \code{graph\_query} (navegação por tipo,
\code{id}, texto ou relação) e \code{timeline\_query} (eventos datados, filtráveis por
entidade, janela e tipo). A alternativa óbvia era RAG no prompt, e ela foi rejeitada por
três motivos. Custo fixo em todo turno para um contexto que a maioria dos turnos não usa.
Perda de rastreabilidade --- contexto colado no prompt é indistinguível de alucinação na
hora em que o cliente pergunta ``como você sabe disso?''. E, principalmente, similaridade
vetorial responde mal a perguntas relacionais e temporais, que são o formato da maioria dos
pedidos aqui: ``o que mudou desde a última reunião'', ``quem aprovou esse criativo''. Aqui
entra o conceito que o enunciado só insinua ao dizer ``linha do tempo'': a
\textbf{bi-temporalidade} das arestas do Graphiti/Zep, que separa \emph{quando o fato passou
a valer} de \emph{quando o sistema soube dele}~\cite{rasmussen2025zep}. É a diferença entre
``o CPA subiu depois que trocamos o criativo'' e ``subiu antes, só descobrimos depois''.
Numa investigação de anomalia, essa distinção \emph{é} a resposta.

\textbf{App} é uma metodologia empacotada que devolve \emph{estrutura}, não parágrafo:
\code{app\_diagnostico} retorna veredito, causas-raiz com evidência e fonte, hipóteses
descartadas e próximos passos. É o que impede a metodologia da agência de virar improviso do
modelo, e é a diferença entre um adaptador opinativo e uma API crua --- a tool certa não é
\code{meta\_ads.call(endpoint, params)}, é
\code{meta\_ads.gasto\_por\_criativo(conta, periodo)}, cuja definição de ``gasto'' é da
agência e não do modelo~\cite{vo2026harness}.

\textbf{Tool}, por fim, é o contrato. A Tabela~\ref{tab:tools} lista o catálogo desenhado:
dez leituras e duas escritas, cada uma amarrada aos nós de onde pode ser chamada. Duas
travas independentes protegem a allowlist --- o orçamento do nó e o campo de nós permitidos
da própria tool --- de modo que um erro de configuração de nó não abra acesso à pausa de
anúncios. Note que \code{propose\_ctas} é leitura: ela produz texto novo, mas não altera
nada fora do harness. Confundir ``gerar'' com ``publicar'' é exatamente como um harness
passa a agir sem gate.


\subsection{O que eu acreditava --- e o que mudou}

No início eu tratava \emph{harness} como sinônimo de \emph{loop}. A pergunta que eu achava
central era: ReAct, grafo, CodeAct ou plan-and-execute? E eu assumia que, escolhido o loop,
a qualidade do agente seria basicamente a qualidade do modelo, com o harness fazendo o papel
de encanamento. Duas coisas mudaram, e uma terceira menor.

\textbf{Primeira: o loop é a parte menos diferenciada do problema.} Todos os cinco tipos
resolvem as quatro tarefas do gestor de alguma forma. O que separa uma resposta que o gestor
consegue levar para a call de uma que ele não consegue defender não é o formato do loop --- é
o \textbf{contrato de fronteira}: o que atravessa para o contexto do modelo, com que rastro,
e o que é permitido em cada fase. Por isso o artefato mais trabalhoso desta arquitetura
acabou sendo a definição de tipos --- o evento de trace como união discriminada
serializável, o preview de ação escrito em português, a entidade resolvida carregando um
grau de confiança --- e não o executor do grafo. Esses tipos \emph{são} o harness. O loop é o
resto.

\textbf{Segunda, e foi um erro concreto que eu cometi no primeiro rascunho: eu ia
implementar permissão como middleware em volta da execução de tool.} Um wrapper que checa o
efeito e decide se deixa passar. Parece limpo, e é o que a maioria dos exemplos faz. Não
funciona neste domínio, por um motivo específico: um middleware pode \emph{negar}, mas não
pode \textbf{parar o turno, montar um preview legível e devolver o controle ao humano com o
trabalho já feito preservado para retomar depois}. Ele só sabe dizer não. E permissão que só
sabe negar empurra o produto para um dos dois extremos ruins: ou o agente vira somente
leitura, ou alguém desliga a checagem. Por isso o \code{gate} virou nó de primeira classe do
grafo, com o turno fechando em ``aguardando confirmação'' e a retomada acontecendo por
reentrada com estado restaurado. A permissão deixou de ser uma checagem e passou a ser uma
\textbf{fase do fluxo de trabalho}. Foi a mudança que mais alterou o desenho --- e ela só
ficou completa quando cruzei com o aviso de idempotência da Seção~3.4, que obriga a separar
o gate do efeito.

\textbf{Terceira, menor mas honesta: eu superestimava RAG.} Achava que despejar o contexto
da conta no prompt de sistema era a forma pragmática de dar memória ao agente. Ao escrever
o passo a passo do diagnóstico, percebi que o valor não está em \emph{ter} o contexto ---
está em conseguir dizer de onde ele veio.

\subsection{Trade-offs aceitos de propósito}

\textbf{Latência por auditabilidade.} O caminho feliz tem, por estimativa de projeto, de
quatro a seis chamadas de LLM, contra duas ou três de um ReAct puro. São segundos a mais em
toda pergunta, inclusive nas triviais: o grafo cobra pedágio até em ``me explica esse
relatório''.

\textbf{Fricção por segurança.} O gate interrompe o turno em toda escrita. Pausar oito
criativos em oito confirmações é pior que fazer no gerenciador; agrupar em uma reduz a
fricção e aumenta a aprovação cega. O trade-off não desaparece, só muda de lado.

\textbf{Cauda longa por superfície pequena.} Sem CodeAct, ``refaz com média móvel de sete
dias'' não tem tool e não tem como ser improvisado. Toda capacidade nova é um ciclo de
desenvolvimento, não uma frase do gestor. Em troca, a resposta a ``o que esse agente
consegue fazer?'' cabe na Tabela~\ref{tab:tools}.

\textbf{Sessão longa por simplicidade.} A compactação é por nó, não por sessão: o contexto
do chat cresce sem teto. É problema conhecido e não resolvido. Pior, \code{compact} é lossy
por definição --- numa conversa de trinta turnos, detalhes antigos somem e o agente não sabe
que sumiram.

\textbf{Fidelidade aritmética por segurança.} Agregação feita pelo LLM erra em silêncio
quando as linhas passam de algumas dezenas. A resposta aqui é fazer a aritmética em código
dentro da tool (Seção~4.5), o que resolve os cruzamentos previstos --- e deixa de pé o custo
direto de recusar o sandbox: o cruzamento que ninguém previu não tem tool. É o assunto da
Seção~6.1.

\textbf{Rigor aparente.} Um grafo bonito dá a impressão de que o sistema é determinístico.
Ele não é: a decisão de qual edge tomar continua sendo do modelo na maioria dos casos.

% ---------------------------------------------------------------------------
\section{Sistema}

A arquitetura das seções anteriores foi implementada e está no ar. A demo é pública em
\url{https://desafio-adz-harness.vercel.app}: responde HTTP 200 sem login e sem cookie,
hospedada na Vercel. Esta seção afirma apenas o que se confirma abrindo essa URL --- todo
número abaixo veio de uma execução contra a rota real do chat, não de estimativa. O que
ficou no desenho está nomeado no fim da seção.

\subsection{O grafo, com os números}

Os sete nós da Figura~\ref{fig:grafo} existem no runtime, com os tetos que a Seção~3.3
descreve em prosa. \code{interpret} roda no máximo três passos, sob allowlist própria de
duas tools --- \code{graph\_query} e \code{timeline\_query} ---, porque resolver entidade
não precisa de API de anúncio. \code{fetch} tem teto de seis passos, \code{reason} de
quatro, e o ciclo \code{fetch} $\leftrightarrow$ \code{reason} fecha em três idas e voltas.
\code{act} tem teto de dois. O catálogo é o da Tabela~\ref{tab:tools}: doze tools, dez de
leitura e duas de escrita, cada uma declarando o efeito e os nós de onde pode ser chamada.

A Tabela~\ref{tab:percursos} mostra os quatro pedidos típicos do gestor medidos nessa rota.
Duas coisas nela interessam mais que as contagens. A primeira: os caminhos \emph{diferem} ---
T1 termina em \code{gate}, com o turno em \code{awaiting\_confirmation}, e os outros três
terminam em \code{respond}. A segunda: T1 entrega os artefatos ao Palco \textbf{antes} da
confirmação. O gestor vê a lista de criativos e os três \emph{diffs} de CTA propostos, e só
então decide. Preview que chega depois da decisão não é preview.

\begin{table}[t]
\centering
\footnotesize
\caption{Os quatro percursos, medidos na rota real do chat. T1 é o único que para no
\code{gate}; os outros três seguem para \code{respond}. A coluna ``Palco'' lista os
artefatos entregues à interface.}
\label{tab:percursos}
\setlength{\tabcolsep}{4pt}
\begin{tabular}{@{}llll@{}}
\toprule
Tarefa & Tools & Termina em & Palco \\
\midrule
T1 pausar criativos & 5 & \code{gate} (halt) & \code{creative\_list}, \\
                    &   &                    & 3$\times$\code{cta\_diff} \\
T2 diagnóstico      & 7 & \code{respond}     & \code{diagnostic}, \\
                    &   &                    & 2$\times$\code{metrics\_table} \\
T3 pauta            & 5 & \code{respond}     & \code{agenda} \\
T4 cruzamento       & 4 & \code{respond}     & \code{metrics\_table} \\
\bottomrule
\end{tabular}
\end{table}

\subsection{Um turno completo, nó por nó}

T2 --- ``por que caíram as vendas da Ômega 3 essa semana?'' --- é o percurso mais longo:
\code{interpret} $\rightarrow$ \code{plan} $\rightarrow$ \code{fetch} $\rightarrow$
\code{reason} $\rightarrow$ \code{respond}, sete tools, 71 quadros NDJSON no stream.

Em \code{interpret}, \code{graph\_query} resolve ``a Ômega 3'' contra o supercérebro e fixa
a janela de datas. \code{plan} não chama tool: ordena os passos. \code{fetch} puxa
\code{meta\_ads\_insights}, \code{crm\_leads}, \code{ga\_report}, \code{list\_criativos} e
\code{timeline\_query} --- e é aqui que aparece o passo que ninguém programou, porque olhar
o link de destino do criativo só faz sentido depois de ver que o Meta reporta mais conversão
do que o CRM atribui. \code{reason} chama \code{app\_diagnostico}, o App de metodologia, que
devolve estrutura: veredito, causas com evidência, hipóteses descartadas. \code{respond}
redige.

O veredito é uma frase: ``As vendas não caíram na proporção que o relatório mostra: a
atribuição quebrou.'' A evidência que a sustenta vem numerada e com fonte: ``Meta reporta 51
conversões de Ômega 3 na semana atual contra 51 em S2; leads atribuídos caíram de 54 para
26, e os leads sem origem subiram de 2 para 30. Receita ganha do produto: R\$ 4.980,33 (S2)
$\rightarrow$ R\$ 5.497,18 (S4).'' O Analytics corrobora pelo lado do canal: ``paid\_social
2939 $\rightarrow$ 2054 sessões; direct 646 $\rightarrow$ 1430'' --- tráfego que existia
deixou de se identificar. Uma causa secundária real, e menor, também sai nomeada: o
carrossel saturou, frequência 5,56, CTR de 1,91\% para 0,72\%. Cada afirmação da resposta
carrega evidência e arquivo-fonte citado, e a lista de hipóteses testadas e descartadas vem
junto --- é a diferença entre um parágrafo correto e um parágrafo defensável na call.

\subsection{O gate é real, e é idempotente}

A Seção~3.4 argumenta que \code{gate} e \code{act} precisam ser nós separados porque a
retomada roda o nó de novo desde o início. Isso deixou de ser teoria. Na rota real, com o
turno T1 parado em \code{awaiting\_confirmation}, três coisas acontecem:

\begin{itemize}\itemsep2pt
\item \textbf{Aprovar} executa a ação \textbf{uma vez}: o grafo entra em \code{act}, emite um
único evento \code{action\_executed} e segue por \code{acao\_confirmada} para \code{respond}.
\item \textbf{Reenviar a mesma aprovação} executa \textbf{zero} vezes. A resposta é ``Não
encontrei um turno aguardando confirmação nesta sessão.'' A pendência já foi consumida.
\item \textbf{Negar} vai direto para \code{respond}, com zero ações executadas e a pendência
zerada.
\end{itemize}

É o elo mais forte entre a tese e o artefato: a regra não-negociável da Seção~3.4 é
verificável clicando duas vezes em confirmar e contando os efeitos.

\subsection{O agente não cai na armadilha do próprio dataset}

O conjunto de dados esconde um caso desenhado para punir quem olha uma fonte só. O anúncio
\code{omega3\_vid\_prova\_social\_v2} tem CPA péssimo no CRM --- e tem porque o rastreio dele
quebrou, não porque parou de vender. Quem lê apenas o CRM desliga o anúncio que está
vendendo.

Em T1, os três anúncios propostos para pausa \textbf{não incluem} esse. Ele aparece na lista
marcado como ``não pausar'', com o motivo. Em T2, ``pausar o prova social por CPA ruim no
CRM'' consta explicitamente entre as hipóteses \emph{descartadas}. É a melhor evidência
disponível de que cruzar fontes muda a decisão, e não apenas enfeita a resposta.

\subsection{Aritmética em código, e o rodapé que declara o que falta}

Todo join Meta $\times$ CRM por \code{utm\_content}, todo agrupamento e todo CPA são
calculados em TypeScript e devolvidos prontos ao modelo, junto com a contagem explícita do
que ficou de fora. O modelo redige sobre um número que ele não somou.

A parte que importa é o resto da frase. O rodapé da tabela de T4 diz, literalmente: ``37
lead(s) do CRM chegaram sem \code{utm\_content} e NÃO entram nesta tabela (R\$ 3.616,93 de
receita ganha). R\$ 900,00 de gasto do Meta e 30 conversões também ficam de fora, do(s)
anúncio(s) \code{omega3\_vid\_prova\_social\_v2} --- a URL parou de carregar UTM. CPL = gasto
$\div$ leads atribuídos: para um anúncio com rastreio quebrado ele fica artificialmente
alto.'' A tabela declara o que a tabela não mostra, e explica por que a métrica dela mente
naquela linha.

Isso mitiga a falha nomeada na Seção~6.1 --- mas mitiga, não elimina, e a distinção é o
ponto. A contagem do que ficou de fora é calculada em código; \emph{exibi-la} depende de o
nó \code{respond} preencher o campo de rodapé. Uma garantia estrutural seria o artefato
recusar-se a existir sem ela. Não é o caso. É uma prática, não uma trava.

\paragraph{O que ficou no desenho.} A compactação existe como nó, mas não há checkpoint
persistente entre sessões: o estado vive no turno. Retry com backoff exponencial, avaliação
quantitativa (item 1 da Seção~7) e qualquer integração real ficaram de fora --- as sete
fontes de dados são arquivos locais, e é disso que trata a seção seguinte.

% ---------------------------------------------------------------------------
\section{Notas de execução}

\subsection{Chave e modelo}

A chave do OpenRouter é colada num campo visível na coluna do chat. Ela fica no
\code{sessionStorage} do browser e viaja apenas no header \code{x-openrouter-key} da chamada
para \code{/api/chat}. Não é lida de variável de ambiente, não é persistida no servidor e
não é gravada em log --- verificado com uma chave-canário, que não reaparece nem no stream
nem no log do servidor. Ao lado do campo, um seletor de modelo oferece
\code{anthropic/claude-sonnet-4.5}, \code{openai/gpt-4.1} e \code{google/gemini-2.5-pro},
mais um campo livre para qualquer outro \emph{slug} do OpenRouter.

\subsection{Sem chave: replay determinístico}

Sem chave colada, a demo entra em modo \emph{replay}: um roteiro gravado dos quatro
percursos, sem nenhuma chamada de LLM. A interface rotula esse modo de forma permanente e
inequívoca, e o rótulo existe por um motivo de honestidade, não de estética --- uma demo que
finge pensar é pior que uma demo que não pensa. O replay serve para ver a topologia, o
trace e o gate; não serve para julgar o modelo.

\subsection{Os dados}

São 544\,KB de JSON em sete arquivos: o supercérebro com 39 nós e 76 arestas, a linha do
tempo com 18 eventos, o Meta com 304 linhas diárias, 397 leads de CRM, 13 criativos, mais
Google Analytics e Google Ads. \textbf{Tudo é fictício.} Nenhum dado de conta real da SPOT
ou da Housewhey foi usado; os nomes vêm do enunciado do desafio, os números foram gerados.

O gerador é determinístico, com \emph{seed} fixa, e vem com um validador de oito checagens
aritméticas: o gasto por anúncio soma o gasto da campanha; todo \code{utm\_content} que
aparece no CRM existe na mídia; leads batem com conversões; o buraco de atribuição da semana
atual supera o da semana anterior; a integridade referencial do grafo fecha. Sem isso, o
caso plantado da Seção~4.4 poderia ser um artefato de geração em vez de um teste.

Registro uma limitação de contrato: o \code{dataset\_prompt.md} oficial não estava mais
acessível na página do desafio na data de execução, então o esquema seguido foi derivado do
próprio enunciado.

\subsection{O que dá para testar}

Onze checagens executáveis, todas passando, cobrem exatamente as afirmações que este paper
faz sobre o harness: escrita fora do nó \code{act} é recusada por política; os tetos de
passo e de ciclo do ReAct valem; o gate é idempotente; a contagem do que fica fora do join
confere; a chave não vaza; e o roteamento observado é idêntico ao da
Figura~\ref{fig:grafo}. Um teste que apenas confirmasse que o sistema responde não provaria
nada --- estes confirmam as fronteiras.

\subsection{O limite honesto}

O caminho com LLM real é exercitado por teste com a chamada de rede \emph{stubada} e pela
rota em produção, mas \textbf{não foi executado ponta a ponta com uma chave paga} --- quem
cola a chave é o avaliador. Portanto este paper não afirma nada sobre a qualidade das
respostas com um modelo real: o que está demonstrado é o comportamento do harness --- a
topologia, os tetos, a allowlist, o gate e a aritmética ---, que é justamente a camada que a
Seção~1 separou do modelo.

% ---------------------------------------------------------------------------
\section{Onde a solução quebra}

Uma falha por tarefa típica do gestor. São falhas do harness que eu propus, não do estado da
arte.

\subsection{Relatório --- a soma silenciosa}

O cruzamento de gasto do Meta com leads do CRM por \code{utm\_content} é feito em código
dentro da tool, e essa é a defesa contra o modo de falha clássico: LLM lendo centenas de
linhas erra soma e agrupamento \emph{em silêncio}, porque o número errado é plausível. A
falha que sobra é outra, e é mais sutil.

Parte dos leads chega sem \code{utm\_content} e simplesmente não aparece na tabela. A
contagem desses órfãos é calculada corretamente --- mas o que impede o artefato de mentir
com cara de rigor é \emph{exibi-la}, e exibi-la depende de o nó \code{respond} preencher um
campo de rodapé. Nada no tipo do artefato o obriga: uma tabela sem rodapé é um artefato
válido. Ou seja, a correção da peça mais importante do relatório continua sendo uma prática,
não uma trava. E permanece o buraco de cobertura: agregação que não tem tool --- ``refaz com
média móvel de sete dias'' --- volta para o colo do modelo, com o mesmo erro silencioso de
sempre.

\subsection{Diagnóstico --- o agente não sabe o que não pode ver}

O harness chega à causa-raiz do diagnóstico por um motivo frágil: uma das tools de leitura
expõe o link de destino do criativo. Se a pista morasse num campo que nenhuma tool expõe ---
uma configuração de rastreamento, um redirect, uma mudança feita direto no gerenciador ---
o agente \textbf{não hesitaria}. Ele concluiria ``queda de demanda na semana'' com a mesma
estrutura de evidência e a mesma confiança, porque ausência de dado não gera sinal. Some-se
o teto de ciclos: uma investigação legítima que precisasse de um ciclo a mais é cortada, e o
corte por orçamento é, para o gestor, indistinguível de ``investiguei e não achei mais
nada''. O sinal de orçamento esgotado existe no estado; ele só ajuda se a resposta o
mencionar.

\subsection{Pauta --- completa na aparência}

A consulta à linha do tempo só conhece o que foi registrado. Na operação real, a decisão
mais importante da semana costuma ter acontecido numa ligação, num áudio de WhatsApp não
transcrito, ou numa mudança que a Aline fez direto no gerenciador e comentou com a Carolina
no corredor. A pauta sai bem formatada, com blocos, responsáveis e prioridades --- e omite
exatamente o item que fará a reunião existir. É o modo de falha mais perigoso dos quatro,
porque a pauta \emph{parece} completa: o gestor entra na call confiando nela em vez de
conferir. Uma pauta vazia seria menos danosa que uma pauta plausível e furada.

\subsection{Criativos --- o gate vira teatro, o apelido não resolve}

Duas falhas na mesma tarefa. \textbf{(a)} O gate depende de o gestor ler o preview. Depois da
décima confirmação, ele clica ``confirmar'' por reflexo, e o preview passa a ser um registro
de que a informação \emph{estava} disponível --- não uma decisão informada. Agrupar
confirmações reduz a fricção e aumenta a aprovação cega; manter uma por item faz o gestor
preferir o gerenciador. Não há saída técnica: o trade-off só muda de lado. \textbf{(b)} A
resolução de entidade quebra com apelido interno. O time chama a campanha de ``a nova da
Ômega''; esse alias não está no grafo; a confiança cai abaixo do limiar e o agente
\emph{pergunta}. Comportamento correto e irritante --- o gestor sabe perfeitamente do que
está falando e o agente parece obtuso. A correção real não é técnica: é curadoria contínua
de aliases no supercérebro, trabalho que o harness não faz sozinho.

% ---------------------------------------------------------------------------
\section{O que eu construiria em seguida}

\paragraph{Avaliação antes de qualquer feature nova.} Um conjunto de 20 a 30 prompts com
resultado esperado e um runner determinístico medindo quatro coisas: a entidade foi
resolvida corretamente; as tools chamadas foram as necessárias \emph{e apenas elas}; a
causa-raiz foi encontrada quando existia; o gate foi acionado em toda escrita. Hoje os tetos
de passos são números que eu estimei olhando quatro prompts --- sem medição, todo ajuste é
chute e toda regressão passa despercebida. É a primeira coisa porque as outras duas ficam
impossíveis de validar sem ela.

\paragraph{Estender a agregação determinística --- e travar o rodapé.} O mecanismo já
existe: o join por \code{utm\_content}, o agrupamento e o custo por lead são aritmética em
código dentro da tool, com a contagem do que ficou de fora (Seção~4.5). O que falta é
cobertura e obrigação. Cobertura: cada recorte novo que o gestor pedir ainda é uma tool
nova, e essa rigidez é o preço aceito por recusar o sandbox --- o código é meu, não gerado.
Obrigação: fazer o tipo do artefato exigir a declaração de exclusões, para que uma tabela
sem rodapé não seja sequer construível. Hoje o rodapé é prática; a Seção~6.1 explica por que
isso não basta.

\paragraph{Aprendizado de alias.} Quando o \code{interpret} pergunta e o gestor corrige
(``é a Ômega 3 Prospecção''), gravar esse alias no supercérebro. Ataca a falha 6.4(b) no
ponto certo: o problema é de curadoria, e a correção mais barata é aproveitar a correção que
o humano já está fazendo de graça.

\subsection*{E o que eu deliberadamente não construiria}

\textbf{Sandbox/CodeAct.} É a tentação óbvia depois do item anterior, e por isso precisa
estar escrita aqui. Código gerado com acesso a credenciais de conta de anúncios de cliente
exige isolamento de verdade --- VM, rede fechada, quotas, revisão de saída --- e nada disso
se constrói bem em uma semana. Meio sandbox é pior que nenhum.

\textbf{Multi-agente com supervisor.} Um agente para Meta, outro para CRM, outro para
criativos. Dobra o custo de tokens, multiplica os pontos de falha e reintroduz um grafo ---
só que implícito, mal especificado e mais difícil de auditar. O grafo explícito já faz esse
trabalho melhor.

\textbf{Escrita autônoma no supercérebro.} O agente inferir fatos da conversa e gravar na
memória sem confirmação. É a coisa mais perigosa da lista: memória é escrita, e memória
errada é pior que dado errado, porque contamina silenciosamente todos os turnos futuros ---
o erro deixa de ter data e vira contexto.

\textbf{Mais conectores de canal.} É a feature mais fácil de vender e a que menos prova a
tese. Cada conector novo é largura sem profundidade.

\textbf{Um modo de autonomia configurável} que permita pular o gate para ações marcadas como
seguras. É exatamente o mecanismo que transforma o gate em teatro por decisão de produto, em
vez de por fadiga do usuário. Se um dia existir, vem depois da telemetria que mostre quais
confirmações são realmente ruído --- não antes.

% ---------------------------------------------------------------------------
\section{Conclusão}

O trabalho de arquitetura de harness, neste domínio, não é escolher entre cinco opções: é
decidir o que fica por fora, o que fica por dentro e o que atravessa tudo. Lastro responde
grafo, ReAct e permissão deny-first, e recusa por escrito o sandbox e o RLM completo, com o
custo de cada recusa nomeado. A aposta que amarra as três decisões é sobre o usuário, não
sobre a tecnologia: o gestor de marketing não precisa de um agente que decida sozinho ---
precisa de um agente cujas conclusões ele consiga defender na frente do cliente, e cujas
ações ele consiga autorizar sabendo o que está autorizando.

\bibliographystyle{plain}
\bibliography{refs}

\end{document}
